% Scope: Derive EPG from spatial Fourier coefficients and the declared rotation convention; track transverse and longitudinal states through RF, relaxation and gradients.
% Status: architecture only; no chapter prose or completed problems yet.
% Prerequisites: Chapters 3 through 8; Chapter 10 for diffusion attenuation.
% Planned worked examples: Stimulated echo, CPMG echo train and imperfect spoiling.
% Planned figures: Signed-order EPG state evolution and pathway diagrams.
% Planned challenge problems: Derive the RF transition matrix; verify EPG against isochromats.
\chapter{Extended Phase Graphs and Signal Pathways}
\label{ch:13}

\section{Configuration-state representation}
\label{sec:13:configuration-state-representation}
\subsection{Spatial harmonics and averaging}
\subsection{F-plus, F-minus and longitudinal Z states}
\subsection{Signed orders, conjugacy and observable echoes}

\section{EPG operators}
\label{sec:13:epg-operators}
\subsection{RF transition matrix from Bloch rotations}
\subsection{Relaxation, recovery and off-resonance}
\subsection{Gradient shifts and state truncation}

\section{Echo pathways}
\label{sec:13:echo-pathways}
\subsection{Hahn echoes and stimulated echoes}
\subsection{Crushers, spoiling and coherence selection}
\subsection{Phase cycling and pathway interference}

\section{Echo trains and steady states}
\label{sec:13:echo-trains-and-steady-states}
\subsection{CPMG and variable-angle FSE/TSE}
\subsection{Spoiled GRE and incomplete spoiling}
\subsection{SSFP and transient MRF applications}

\section{Extensions and numerical verification}
\label{sec:13:extensions-and-numerical-verification}
\subsection{Diffusion weighting along pathways}
\subsection{Slice profiles and noninteger dephasing}
\subsection{Convergence and comparison with Bloch ensembles}

\section{Worked examples}
\label{sec:13:worked-examples}
% Planned calculations: Stimulated echo, CPMG echo train and imperfect spoiling.
\subsection{Derivation and quantitative calculation}
\subsection{Assumptions, limiting cases and scanner interpretation}

\section{Challenge problems}
\label{sec:13:challenge-problems}
% Problem design: Derive the RF transition matrix; verify EPG against isochromats.
% Use challengeproblem and stable labels prob:13:<topic>.
% Complete solutions belong in Appendix E, section for this chapter.
% Use \begin{solution}{prob:13:<topic>} ... \end{solution} there.
\subsection{Derivation and quantitative design}
\subsection{Model criticism and integrated reasoning}
